\section{Introduction}

% Program rewriting is at the core of many modern program optimizers~\citep{halide, TODO}.

Equality saturation (\eqsat) and Datalog are both
 fixpoint reasoning frameworks with 
 many 
 applications,
 extensions,
 and high-quality implementations~\cite{egg, souffle}.
They share a common setup: 
 the user provides \emph{rules} and an initial set of \emph{facts} 
   (a term in \eqsat and a database in Datalog),
 then the system derives a larger and larger set of facts from those inputs.
However,
 their commonalities have not---until now---been fully realized or exploited.
As a result, 
 the frameworks have developed independently and are used in different domains.
Datalog is well-studied by the databases community, 
 and practitioners use modern implementations
 to build program analyses% 
 ~\cite{cclyzer,cclyzerpp,doop-datalog}.
%  several implementations also target commercial applications.
Equality saturation
 is a more recent, term-centric technique
 favored in the programming languages community
 for program optimization and verification.

As users apply EqSat and Datalog to new, 
  more demanding problems,
  the limitations of each tool
  become apparent.
For example,
 Herbie~\cite{herbie}, 
 a tool that uses \eqsat to optimize floating-point accuracy, 
 relies on unsound rewrites because it lacks the analyses
 to prove that certain rewrites are safe 
 (e.g. $x/x \rightarrow 1$ only if $x \neq 0$).
To combat the unsoundness,
 Herbie must validate the results of \eqsat
 and discard them if unsoundness was detected.
On the Datalog side,
 \cclyzerpp~\cite{cclyzerpp},
 a recent points-to analysis system implemented in Datalog
 that supports Steensgaard analyses~\citep{DBLP:conf/popl/Steensgaard96}
 for LLVM~\citep{llvm}
 resorted to an ad-hoc implementation of union-find,
 because the provided implementation of equivalence relations
 was too slow.
The resulting implementation's complexity led to bugs
 in the pointer analysis.
In short, EqSat struggles to support rich analyses, 
  and equational reasoning in Datalog is complex and slow.

\emph{Our key insight is that
 the efficient equational reasoning of EqSat and 
 the rich, composable semantic analyses of Datalog 
  make up for each other's weaknesses, 
  and unifying the two paradigms brings 
  together---and goes beyond---the best of both worlds.}
In fact, spontaneous developments in both communities 
  have already begun converging towards each other:
Datalog tools have 
  added efficient equivalence relations~\cite{eqrel},
  lattices~\cite{flix,ascent},
  and some support for datatypes~\cite{souffle-adt}, 
  while the \eqsat community
  has recently developed support for
  conditional rewriting,
  lattice-based analyses~\cite{egg,metatheory.jl},
  and relational pattern matching~\cite{relational-ematching}.
We bring this trend to completion and close the gap
  between EqSat and Datalog.

In this work,
 we propose \egglog,
 a fixpoint reasoning system that subsumes both \eqsat and Datalog.
It contains all of the innovations listed above as well as new ones,
 and it addresses crucial limitations that have prevented progress in 
 real-world applications.
\egglog is essentially a Datalog engine
 with two main extensions. % that have far-reaching consequences.
First,
 \egglog has a built-in, extensible notion of equality.
The user can assert that two terms are equivalent,
 from which point on they are indistinguishable to the system.
For example,
 consider a relation with a single tuple: $R = \set{(a, b)}$
 for distinct $a$ and $b$.
The query $R(x, x)$ would yield nothing,
 but if the user asserts that $a$ and $b$ are equivalent,
 then the query would return the equivalence class containing both $a$ and $b$.
% \mw{also unification variables}
Second,
 \egglog has built-in support for (uninterpreted) functions.
From a relational perspective,
 a function is a relation with a functional dependency 
 from its arguments to its output,
 i.e., the output is uniquely determined by the arguments.
However, user-extensible equality
 introduces challenges for maintaining functional dependencies.
Consider a function $f$ such that $f(a) = b$ and $f(c) = d$, 
 but $b$ does not (yet) equal $d$.
What happens when the user asserts that $a$ and $c$ are equivalent?
An \egglog function can be annotated with a \emph{merge expression},
 a novel mechanism that 
 \egglog uses to resolve functional dependency violations
 by combining the two conflicting output values.
In the above case,
 $f$'s merge expression might 
 assert that $b = d$ (essentially asserting congruence of $f$),
 or return the supremum of $b$ and $d$.
The flexibility of merge expressions allows \egglog
 to exceed the expressive power of both \eqsat and 
 Datalog extensions with lattices.
% \mw{following para seems redundant, jrw agrees}
% \rw{The paragraph maps the solutions to the problems.
%  Without it it's unclear how the extensions have anything to do with
%   analyses and rewriting.  }
% \Egglog allows users to seamlessly combine 
%   analyses and equational reasoning. 
% Each function can map a term to either an equivalence class 
%   or a lattice value in some abstract domain.
% Then, the merge expression can either unify equivalence classes 
%   or take the join of lattice values.
The high-level \Egglog language allows the user to specify 
  complex interactions among terms, equivalence classes, 
  and lattice values.
At the same time, highly optimized algorithms for 
  relational and equational reasoning 
  work together to make \Egglog efficient.

The combination of \eqsat and Datalog
 also brings many practical---if somewhat more prosaic---benefits.
For example,
 \citet{relational-ematching}
 observed that \eqsat is hampered by inefficient
 \ematching (pattern matching modulo equality) algorithms,
 and that a relational approach can be vastly more efficient.
\egglog's Datalog-first design naturally 
 supports efficient \ematching by reducing it to a relational query.
This goes even further:
incremental \ematching is only supported in some SMT solvers like Z3~\cite{z3}
 and has not yet made its way into \eqsat implementations,
 while \egglog supports them \emph{for free} with semi-na\"ive evaluation~\citep{seminaive},
 a common technique that makes Datalog incremental.
\egglog's support for functions
 provides the basis for working with terms,
 which only have limited support in other Datalog systems~\cite{souffle-adt}.
Users can also define multiple
 functions and datatypes
 to model their domain,
 unlike most \eqsat tools~\cite{egg,metatheory.jl}
 where users are forced to use a single, ad-hoc datatype.
Finally,
 \egglog is designed as a language (as well as a library),
 making it more accessible
 than \eqsat libraries~\cite{egg,metatheory.jl} 
 that are locked to their implementation language.
% \yz{
%   The claims made in this paragraph are not substantiated in the rest of the paper.
%   MW: which claims specifically?
% }

We perform two case studies showing that \egglog out-performs 
  state-of-the-art applications of \eqsat and Datalog respectively.
First, we show that \egglog makes Steensgaard-style points-to analyses
  faster and easier to write.
Compared to the \souffle Datalog system,\xspace \egglog computes the points-to analysis
  4.96$\times$ faster.
Second, we demonstrate the power of \egglog
 with a new, sound implementation of Herbie's \eqsat procedure.
This allows Herbie to perform aggressive optimizations soundly,
  and return results faster given the same error tolerance.

% In both cases,
%  the \egglog implementation is simpler, faster, 
%  and corrects bugs from the original implementation.
% \of{We don't have evidence we are faster than Herbie yet! We might by tonight}

In summary, this paper makes the following contributions:
\begin{compactitem}
  \item We introduce a bottom-up, Datalog-like logic language 
   for equality saturation and similar unification-based algorithms.
  \item We present a fixpoint semantics for the core language of \egglog.
  \item We present an implementation for \egglog with optimizations 
   from database research such as semi-na\"ive evaluation.
\end{compactitem}

% Composing analysis easily using datalog
% Making analysis more powerful with the power of equality
% Performance improvement (seminaive?)





%%%%%%%%%%%%%%%%%%%%%%%%%%%%%%%%%%%
%%%%%%%%% old stuff below %%%%%%%%%
%%%%%%%%%%%%%%%%%%%%%%%%%%%%%%%%%%%



\begin{comment}

\section{Alternative Introduction}

Program optimization and program analysis are 
  two important tasks in programming languages research.
Researchers have developed different theories and tools for each of these tasks.
In program analysis, Datalog has been successfully applied to
  specify analyses in a declarative yet scalable way.
In optimization, equality saturation has seen recent breakthroughs to make it effective in many domains.
In this work, we observe that both paradigms are based on the same principle of 
  repeatedly applying a set of rules.
In other words, they are both \emph{fixpoint} algorithms. 
Building on this common core, 
  we propose the first framework that combines equality saturation 
  and datalog within the same system.
Our system, called egglog, 
  allows the user to specify tightly integrated rewrite and analysis rules 
  using the same high-level language.
We have used egglog to re-implement an optimizer built on equality saturation,
  as well as a program analyzer built on Datalog.
Doing so, we made the optimizer sound 
  by augmenting its rewrite rules with semantic information from analyses.
We also made the analyzer more efficient, thanks to
  fast congruence closure algorithms developed for equality saturation. 

A Datalog-based program analyzer works over a program represented as \emph{relations},  
  where each relation is a set of \emph{facts}. 
A fact conveys structural information about the program, 
  from the source code or a more detailed representation such as the control flow graph.
For example, the fact \lstinline|store(x, y)| is generated from the assignment \lstinline|*x = y|.
Given a set of such structural facts, the analyzer infers further \emph{semantic} facts 
  by applying a set of inference rules to the facts.
For example, the rule \lstinline|vpt(x, y) :- store(x, y)| says that
  if there is a statement assigning the value of \lstinline|y| to the address \lstinline|x|,
  then \lstinline|x| is a pointer to \lstinline|y|.
An inference rule can also be recursive,
  for example the rule \lstinline|vpt(x, y) :- vpt(x, z), vpt(z, y)|
  says \lstinline|vpt| is transitively closed.
The analyzer applies the rules to infer new facts, 
  and then applies the rules on the union of old and new facts to infer more facts.
This process continues until no new facts can be found, 
  and the analyzer outputs the final set of facts.
In our example, the \lstinline|vpt| relation gives a points-to graph for the input program.
Specifying analyses in Datalog has many benefits.
First, the declarative nature of Datalog allows one to specify complex analyses with 
  a small set of rules. \rw{Cite evidence showing Datalog analyses are short (might not be true).}
Second, analyses written in Datalog are \emph{composable}.
Since each analysis is specified with a set of rules,
  and the rules can be arbitrarily mutually recursive,
  one may combine multiple analyses by simply ``hooking up'' their rule sets.
Third, thanks to decades of research optimizing the performance of Datalog engines,
  Datalog analyses are scalable.
\rw{Cite Souffle, DOOP etc. to support scalability.}

An optimizer based on equality saturation also works over a program represented as relations.
A relation in equality saturation typically represents syntactic strucure of the program.
For example, the relation \lstinline|+(c, x, y)| roughly says ``the AST node c has + as the operator, 
  and x and y are its operands''.
One relation is special: the \emph{equivalence} relation groups programs and subprograms 
  into \emph{equivalence classes}, 
  where two programs are in the same equivalence class if they are equivalent.
In the beginning, each (sub)-term is in its own equivalence class.
Then the optimizer expands the classes by applying a set of \emph{rewrite rules}.
For example, the rule \lstinline|(+ x y) => (+ y x)| matches on all \lstinline|+|-nodes
  and rewrites the matched term by commutativity. 
Crucially, a rewrite in equality saturation does not delete the rewritten term, 
  but instead adds a new term and asserts that the two terms are equivalent.
In other words, the rewrite rule above can be broken down into two rules in the style of Datalog:
  the first rule \lstinline|+(c2, y, x) :- +(c1, x, y)| says that if there is a \lstinline|+|-node \lstinline|c|
  with operands \lstinline|x| and \lstinline|y|,
  then we add a new \lstinline|+|-node \lstinline|c'| with operands \lstinline|y| and \lstinline|x|;
  the second rule \lstinline|c1 = c2 :- +(c1, x, y), +(c2, y, x)| says that if there are two \lstinline|+|-nodes
  with their operands swapped, then they are equivalent.

The fact that we can express rewrite rules in the style of Datalog suggests
  that it is possible to implement equality saturation using Datalog.
Doing so brings many benefits. 
First, we can benefit from the powerful optimizations developed for Datalog 
  to speed up equality saturation.
Such optimizations include semi-na\"ive evaluation, 
  magic set transformation, and join ordering.
Second, we can leverage the composability of Datalog analyses to 
  gain finer control over the rewrites.
For example, an interval analysis allows us to deduce a term 
  is never zero, if its interval does not contain zero, 
  then we can specify the (\emph{sound}) division cancellation rule
  \lstinline|(/ x x) => 1 if x != 0|. 
Finally, the long and fruitful research on Datalog 
  has produced many theoretical results, 
  and these results may help us better understand equality saturation.
For example, equality saturation may fail to terminate in surprising ways, 
  but researchers have developed precise termination conditions for 
  various extensions of Datalog\footnote{Termination has been extensively studied in the context of \emph{the Chase} 
  which generalizes both Datalog and equality saturation.}, 
  some of which generalize equality saturation.

Nevertheless, a na\"ive encoding of equality saturation in Datalog can be inefficient.
The simplest encoding uses a vanilla binary relation \lstinline|eq(c1, c2)| to represent
  the equivalence relation,
  and supply the reflexivity, symmetry, and transitivity rules for this relation.
Such a relation takes space quadratic to the size of each equivalence class, 
  whereas the union-find data structure used in standard equality saturation implementations 
  take only linear space.
Although certain Datalog engines like Souffle support representing equivalence relations 
  with union-find data structures, 
  simply using this feature is not enough.
The reason is that, when matching the LHS of a rewrite rule, 
  we must match on terms \emph{modulo equivalence}.
For example, the LHS of the associativity rule \lstinline|(+ x (+ y z)) => (+ (+ x y) z)|
  expands to \lstinline|+(c, x, c1), +(c2, y, z), eq(c1, c2)|.
Running this rule as-is requires two joins among three relations, 
  whereas the standard equality saturation implementation only needs one join.
Aside from efficiency, there are also semantic challenges in supporting equality saturation in Datalog.
An important assumption in Datalog is that the rules are \emph{monotonic},
  which allows the rules to be applied in any order. 
The evaluation algorithm and various optimizations for Datalog 
  were designed under this assumption.
Yet it is not clear if the complex rules required to implement 
  rewriting and analyses in equality saturation are monotonic.
\rw{IDK how much we want to care about monotonicity in this paper. 
I just want to add another challenge outside of efficiency. Feel free to change.}
Having experimented with many Datalog engines, 
  we have found none of them satisfy all of our requirements. 
Therefore, we opted to design and implement a new engine, egglog, 
  to combine the best of Datalog and equality saturation.

\rw{Below is pulled from old text.}
 We propose \egglog, a logic language that supports composable rewrites, analyses, and extraction.
\Egglog combines the power of equality saturation and fixpoint reasoning \textit{a la} Datalog.
To combines these two, which uses two totally different representations (term-based v.s.\ relation-based), 
 we extend previous work on relational encoding of e-graphs with rewrites and congruence maintenance, 
 which we call relational e-graphs.
The resulting framework is very powerful, with applications beyond sound rewrites and precise extractions.

For example, it is known that different kinds of semantic analyses can have beneficial interactions with each other~\citep{TODO}:
a constant propagation analysis can be used to determine 
 that certain branching conditions are equivalent to boolean constants,
 which allows dead-code elimination to eliminate more branches,
 which may in turn help the constant propagation.
While manually composing analyses requires careful design of interactions,
 \citet{lerner-analyses-transformations}
 shows that program rewriting can be used
 as a common ground for automatically communicating analyses.
Not only being able to express this composable analyses technique,
 our framework indeed improves it further by incorporating the power of equality saturation,
 which allows even more precise analyses to be derived via non-destructive rewriting.

Moreover, potentially mention sketch-guided eqsat\dots

Compared to the traditional e-graph representation,
 our relational e-graphs also makes equality saturation more efficient.
For example, we are able to apply semi-na\"ive evaluation, a standard optimization in database literature,
 to incrementalize equality saturation.
We are also able to use the recently proposed relational e-matching algorithm~\citep{relational-ematching}:
Relational e-matching views the pattern matching problem in equality saturation as a relational query
 and provides orders-of-magnitude speedups for complex queries.
However, to have both efficient rewriting and efficient matching, 
 a traditional equality saturation framework needs to switch back and forth
 between the e-graph and its relational representation, which can be costly.
\Egglog commits to a single relational representation
 and relationalizes the equality saturation algorithm itself,
 thus eliminating this cost.


%% even older text
 
In fact,
 recently the systems have been developing similar features in parallel.
Both have recently grown features for working with lattice-based domains,
 and 


Equality Saturation engines have recently provided state-of-the-art results
across a diverse range of program optimization, synthesis, and verification tasks.
A key enabler for EqSat is that it is based on congruence and e-graphs.
This means that EqSat search can explore vast spaces of equivalent candidate
terms using only a modest amount of space.
Unfortunately, analyses do not compose in EqSat and just using them requires
writing custom Rust code and learning a new API.

Datalog has been successful for decades.
Recent Datalog extensions (e.g., in Souffle) have enabled state-of-the-art results
in program analysis.
A key enabler for Datalog is that rules compose and Datalog provides a unified
mental model for program execution: relations and rules.
Unfortunately, Datalog engines do not directly support congruence-based reasoning,
and simple efforts to encode it encounter quadratic overhead.
This limits the scale of terms that can be explored with Datalog-only approaches.

In this paper, we propose a new solver abstraction that provides the best
of both EqSat and Datalog's key features: a minimal fixpoint API in which
the programmer can work in a unified, Datalog-like model of relations and queries where
``everything is (composable) rules'', but
which integrates congruence-based reasoning to enable EqSat-scale search.




Program optimizers often stratify rewriting into different phases and apply them sequentially.
However,
 the order of rewrite application
 could drastically affect the quality of optimized programs,
 since applying one rewrite may disable other rewrite from firing.
This is known as the phase-ordering problem in the literature.
Equality saturation~\citep{eqsat} mitigates this problem by
 applying all the rewrites at the same time \textit{non-destructively}.
Equality saturation additionally uses e-graphs to compactly store the large set of generated programs.
The equality saturation approach has been shown to be successful in recent years, 
 with applications in a diverse range of domains including floating-point arithmetic \cite{herbie},
 rule synthesis \cite{ruler}, and query optimization~\citep{wang2022optimizing}.
% New rewrite system implementations
%  are also proposed that are specialized to the equality saturation workloads~\citep{egg}.

Despite the success,
 many equality saturation applications are limited by the expressivity of rewrite rules.
For example,
 Herbie, a system that uses equality saturation to improve the accuracy of floating-point expressions,
 uses unsound rules (e.g., $a/a\rightarrow 1$, which is unsound when $a=0$) 
 to simplify expressions.
To ensure soundness under the presence of unsound rules,
Herbie stops rewriting immediately 
 and returns the best result so far when inconsistencies are discovered,.
Sound rewrites in these cases are difficult if not impossible to express syntactically.
% As a result, semantic analyses are needed (e.g., a program analysis that shows $a$ is never equal to $0$)).
To make the rewrite rule sound, semantic analyses need to prove certain properties over the (set of) terms under rewriting
 (e.g., $a$ is never equal to $0$).
Many equality saturation systems like egg support simple semantic analyses.
In egg, an e-class analysis associates a lattice value to each equivalent class (or e-class) of programs.
When an e-class is modified, changes to e-class analyses are propagated bottom-up.
However, multiple analyses cannot co-exist in a single e-graph, 
 so the semantic analyses is not directly composable in these frameworks.
For example, the developers cannot specify three independent analyses 
 that track the upper bound, lower bound, and non-zeroability of an e-class
 and communicate with each other.

The limited expressiveness also affects the quality of the final extracted programs.
Extractions are often used at the end of equality saturation to extract from a large equivalence class of programs
 subject to given cost functions.
Many applications define the cost functions based on the node types and simple analyses,
 which are nonetheless insufficient in some cases.
For example, an application that optimizes lambda calculus may want to define the cost function based on the type.
This is not doable in frameworks like egg because e-class analysis is computed bottom up, 
 while type checking needs to propagate the context information top down.

To solve these issues, 
 we propose \egglog, a logic language that supports composable rewrites, analyses, and extraction.
\Egglog combines the power of equality saturation and fixpoint reasoning \textit{a la} Datalog.
To combines these two, which uses two totally different representations (term-based v.s.\ relation-based), 
 we extend previous work on relational encoding of e-graphs with rewrites and congruence maintenance, 
 which we call relational e-graphs.
The resulting framework is very powerful, with applications beyond sound rewrites and precise extractions.

For example, it is known that different kinds of semantic analyses can have beneficial interactions with each other~\citep{TODO}:
a constant propagation analysis can be used to determine 
 that certain branching conditions are equivalent to boolean constants,
 which allows dead-code elimination to eliminate more branches,
 which may in turn help the constant propagation.
While manually composing analyses requires careful design of interactions,
 \citet{lerner-analyses-transformations}
 shows that program rewriting can be used
 as a common ground for automatically communicating analyses.
Not only being able to express this composable analyses technique,
 our framework indeed improves it further by incorporating the power of equality saturation,
 which allows even more precise analyses to be derived via non-destructive rewriting.

Moreover, potentially mention sketch-guided eqsat\dots

Compared to the traditional e-graph representation,
 our relational e-graphs also makes equality saturation more efficient.
For example, we are able to apply semi-na\"ive evaluation, a standard optimization in database literature,
 to incrementalize equality saturation.
We are also able to use the recently proposed relational e-matching algorithm~\citep{relational-ematching}:
Relational e-matching views the pattern matching problem in equality saturation as a relational query
 and provides orders-of-magnitude speedups for complex queries.
However, to have both efficient rewriting and efficient matching, 
 a traditional equality saturation framework needs to switch back and forth
 between the e-graph and its relational representation, which can be costly.
\Egglog commits to a single relational representation
 and relationalizes the equality saturation algorithm itself,
 thus eliminating this cost.

% Semantic analyses are not only used during rewriting:
%  in the case of EqSat, they are also used as part of the cost functions when users want 
%   to extract optimal programs from the equivalent classes of programs.
% For example, Storel~\citep{todo}, a system that optimizes tensor storages using EqSat,
%  uses lambda calculus as its intermediate representations, 
%  and the developers wanted to get the type and shape information of each e-class of their tensor programs 
%  with e-class analysis but failed,
%  because the type and shape of a term depends not only on the term itself, but also the context it is in, 
%  which is nonetheless propagated top down.
% yz: maybe we should just not mention multi-patterns in the intro
% For example, Tensat, which uses EqSat for optimizing tensor computation graphs, 
%  rewrite multiple terms together to share computation with rules like
%  \[
%   \begin{cases}
%       m_1\cdot m_2\rightarrow \textit{split}_0(m)\\
%       m_1\cdot m_3\rightarrow \textit{split}_1(m)
%   \end{cases}
%  \]
% where $m$ is $m_1\cdot \textit{concat}(m_2, m_3)$.
% This is known in the literature as multi-patterns.



\end{comment}